\section{Conclusion}
\label{sec:conclusion}

In this paper, we introduced NSE, a Non-Destructive Supervision Extraction framework for Noisy Partial Label Learning (NPLL). NSE shows that high-performance NPLL can be achieved without cross-epoch persistent source write-back. The key is not to leave weak supervision unprocessed, but to restore the native candidate/crowd prior at the beginning of each epoch and extract reliable active supervision from the restored working prior.

Topology-DAES is the core extraction module of NSE. By combining masked-entropy topology reliability with entropy-adaptive KNN propagation, it suppresses ambiguous neighborhoods and strengthens consistent weak-label neighborhoods. Prior-constrained model evidence and active-set training are then built on top of this topology-aware geometric belief. These results suggest that the future direction of NPLL should not be limited to candidate-set purification, reconstruction, or persistent pseudo-target carry-over. A complementary and effective direction is to decouple the native weak-supervision record from the epoch-local working supervision and improve how temporary active supervision is extracted from it.

NSE still depends on the quality of the learned representation and the semantic
purity of local neighborhoods. When neighborhoods are strongly mixed or the
crowd/candidate prior provides little recoverable structure, epoch-wise
restoration alone cannot guarantee reliable extraction. This limitation makes
active-set precision, coverage, and source-drift diagnostics important parts of
evaluating future NPLL methods.

